\chapter{Methodology}
\label{chap:met}

This chapter describes the methodology used to implement language server for the
Jasmin programming language. The primary goal is to build a working language
server prototype and develop the integration for one of the text editors that
support LSP for further extensions and improvements. The following sections
provide the requirements, general server architecture, testing methodology
and other design considerations.

\section{Requirements and Development approach}

In this section we detail the requirements we pose onto the language server
features.

\subsection{Requirements}
\label{chap:met:requierements}

The design of the language server necessitates support for core language
intelligence features, which are outlined below:

\begin{itemize}
  \item Symbol hover functionality (type ascription)
  \item Go to definition
  \item Syntax/semantic highlight support
  \item Source code diagnostics (parsing, typechecking)
\end{itemize}

Also it is important to have not only the language server executable binary
itself, but also clear instructions on how to use the language server in one's
text editor of choice.

\subsection{Design}

Considering the OCaml and Semgrep language server design mentioned in
\ref{sec:background:implementations}, we have chosen to follow the same approach
and implement the language server in a way that our implementation must be
easily extendable and modifiable, allowing  and process each client request
asynchronously.

Judging by the requirements, the language server would be implemented in OCaml, 
since the Jasmin compiler is written in OCaml, allowing easy integration 
with the compiler API. 

The Jasmin language server implementation can be divided into two core
components: JSON-RPC Server and a group of LSP handlers. RPC server represents
the communication bus between the LSP client and the language server. 

Server handles low level protocol operations, including parsing raw JSON-RPC
messages, converting these messages between JSON and OCaml internal
representation, routing the messages to the corresponding handlers and error
handling for malformed messages.

LSP handlers, in turn can be divided in two groups too: the request handlers and
the notification handlers. Request handlers, as expected, process LSP requests,
e.g. \texttt{textDocument/definition}, \texttt{textDocument/hover}. It is worth 
noting that each request must be completed within a feasible timeout, also it is
important to keep in mind during the implementation, that the ordering of the
requests is important too. The server may reorder independent requests like
\texttt{textDocument/completion} and \texttt{textDocument/signatureHelp}, as
they don't affect each other's results. However, dependent requests like
\texttt{textDocument/definition} and \texttt{textDocument/rename} must maintain
order, since renaming could invalidate definition results. Notification handlers, 
in general, do not require immediate server response, in some cases
(\texttt{inialized}) notification does not require the server response at all.
Both handler groups share the access to the compiler's API for semantic analysis,
and the shared server state. 

To implement syntax or semantic highlight support for the code, one might 
implement the corresponding LSP capability, which will collect all semantic 
tokens of the source code tree and then map these tokens to the LSP tokens.  
In our case we are going to implement the tree-sitter for the syntax highlight
instead, getting the syntax highlight support only.

For more information on tree-sitter and its usage in this work please refer to
Section~\ref{impl:tree-sitter}.

\subsection{Development approach}

Instead of releasing the whole language server implementation at once, we chose
to follow iterative development approach, splitting the development process into 
stages. The main development stages outlined below.

\subsubsection*{Implementing syntax highlight support}
First, we implement the syntax highlight support utilizing existing Jasmin
grammar definition. We transform the BNF-like grammar for Menhir~\cite{menhir}
to tree-sitter Javascript, and then test the obtained grammar on the corpus of
correct (in terms of grammar) Jasmin files. These files can be found in the 
libjade~\cite{gh-libjade} repository.

\subsubsection*{Implementing core LS initialization routine}
We implement the JSON-RPC event loop handler, which processes raw
JSON-RPC messages (e.g., \texttt{initialize}, \texttt{textDocument/didOpen}
capabilities).  This includes parsing and validating incoming requests and
responses and managing the Language Server state.

\subsubsection*{Adding diagnostics support}
We integrate the compiler's API to parse the source into the CST, type-check
the CST and transform it into the AST, map compiler diagnostics to the
LSP-compliant diagnostics and support real-time updates, e.g. on file
changes.

\subsubsection*{Adding code intelligence and documentation}
We proceed with further extension of language server capabilities by extending
AST interaction to provide symbol information (e.g., definitions, references).
We implement functionality to get the symbol type on cursor hover, as well as 
write clear documentation on how to use the language server.

\section{Testing methodology}

Since our implementation relies on the compiler-provided parser and typechecker
as an external library, we assume their correctness and do not test their core
functionality.  

For the remaining components, we employ the following testing strategies:

\begin{itemize}
  \item \textbf{Unit Testing.} We verify the correctness of AST transformations,
  code generation, and other critical logic. These tests ensure proper
  interaction with the compiler's API (e.g., correct function calls, error
  handling).  
  
  \item \textbf{Manual Testing} We perform hands-on verification by running the
  compiled binary and observing its behavior with an LSP client (Neovim). 

\end{itemize}
